\section{Introduction}
\label{sec:intro}
% Budget: ~0.75 page.

\begin{outline}
  \item \textbf{Problem.} In team capstones, staff need to know how teams are working, but
    the usual instrument (peer ratings of contribution) only measures contribution, is
    collected at set points, and students can game or soften it
    (cite kaufman2000accounting, ohland2012catme, hooshangi2025assessment).
  \item \textbf{The untapped source.} Students already write regular reflective journals,
    including a prompted section on team dynamics. These hold a running account of how
    each team is going, but the volume is beyond what staff can read closely: one round in
    our course was 222 journals, median about 1,500 words, about 370,000 words in total
    (\S\ref{sec:context}).
  \item \textbf{What we built.} A tool that reads one round of a team's journals with an
    open-weight LLM run on university hardware, and flags three kinds of team trouble
    (conflict, unequal contribution, coordination problems), each backed by verbatim quotes.
  \item \textbf{What this paper reports.} The course context, the tool's design and how it
    evolved, its use on one live round of journals that six tutors reviewed, and a
    reflection on what worked and what did not.
  \item \textbf{Contributions} (short list): (1) a reusable design for LLM journal flagging
    under privacy constraints; (2) evidence on its consistency and on how tutors judged its
    flags; (3) lessons for staff considering something similar.
\end{outline}
